% MIN 27:00 à 28:00. Contributions personnelles, papier compagnon et acquis.
\begin{frame}{Mes contributions au projet}
    Mon travail a porté sur les agrégateurs, \code{Model}, les architectures de modèles et les
    simulations centralisés. J'ai écrit les tests associés automatisés.

    J'ai aussi entièrement rédigé la documentation et les tutoriels. J'ai
    assuré une partie de la rédaction du papier compagnon, puis l'ai corrigé
    avec El-Mahdi El-Mhamdi.

    \smallskip
    {\small Le papier est en cours de rédaction. Il vise la rubrique logiciels
    libres du \textit{Journal of Machine Learning Research} (JMLR).}

    \medskip
    Le stage m'a appris à justifier mes décisions d'architectures et à faire relire mon code.
    J'ai commencé à apprendre à lire des articles de recherche. J'ai aussi
    découvert la rédaction d'un article scientifique.
\end{frame}

% MIN 28:00 à 29:00. Bilan des travaux, bénéfices ISF et prolongement en thèse.
\begin{frame}{Les acquis du stage et la suite}
    Krum réunit des composants réutilisables pour reproduire et comparer
    des expériences d'agrégation robuste.

    \bigskip
    \begin{columns}[T,onlytextwidth]
        \begin{column}{0.47\textwidth}
            \begin{block}{Les acquis ISF en pratique}
                J'ai relié statistiques et optimisation au code testé,
                puis à l'analyse des résultats.

                Le travail en équipe m'a appris à expliquer et documenter mes choix.
            \end{block}
        \end{column}
        \begin{column}{0.47\textwidth}
            \begin{block}{La thèse dans la chaire ATLAS}
                J'étudierai les garanties des agrégateurs pour la sécurité de l'apprentissage.

                Krum servira à mener des campagnes expérimentales comparables.
            \end{block}
        \end{column}
    \end{columns}

    \bigskip
    \centering
    Merci à l'équipe et merci de votre attention.
\end{frame}

% MIN 29:00 à 30:00. Cette minute reste une marge pour l'exposé principal.
